\documentclass[11pt]{article}
\usepackage[margin=1in]{geometry}
\usepackage{setspace}
\usepackage{parskip} % nicer paragraph spacing, no indentation
\setstretch{1.10}

\begin{document}

% ===== Header =====
\begin{flushleft}
\textbf{Zhongli Wang}\\
Bielefeld, Germany\\
Phone: +49 160 9541 6015\\
Email: nkwangzhongli@gmail.com\\
Website: https://sites.google.com/view/zhongli-wang
\end{flushleft}

\vspace{0.3cm}

% ===== Recipient =====
\begin{flushleft}
\textbf{Dr.\ Asel Doranova}\\
Tilburg School of Economics and Management (TiSEM)\\
Tilburg University\\
Tilburg, The Netherlands
\end{flushleft}

\vspace{0.3cm}

\begin{flushright}
\today
\end{flushright}
\vspace{0.3cm}


% ===== Subject =====
\textbf{Application for the Postdoctoral Researcher Position in Economics (HySUCCESS)}\\

\vspace{0.3cm}

Dear Dr. Doranova and Members of the Selection Committee,

I am writing to apply for the \textit{Postdoctoral Researcher in Economics --- HySUCCESS} position at Tilburg University (Department of Economics, jointly with TILEC). I am currently completing my Ph.D.\ in Economics at Bielefeld University and Università Cattolica del Sacro Cuore (Marie Skłodowska--Curie EPOC programme). My research sits at the intersection of industrial organization, market design, and computational economic modeling, with a focus on how market structure, long-term incentives, and policy rules shape investment and competitive outcomes in emerging markets. This makes HySUCCESS an excellent match, particularly given its emphasis on hydrogen market design, long-term contracting, and governance for Europe’s energy transition.

My doctoral research develops dynamic models of strategic firm behavior under infrastructure constraints. In my dissertation, \emph{Strategic Behavior and Economic Dynamics in ICT Markets: An Agent-Based Framework}, firms choose capacity investment and pricing while demand is shaped by network effects, congestion, and learning. Although the setting is digital markets, the economic mechanisms are directly relevant for hydrogen: forward-looking investments, demand and supply uncertainty, and policy-induced changes in incentives. Methodologically, I combine economic theory with simulation-based analysis to evaluate counterfactual designs and to understand how alternative rules affect market performance and welfare.

HySUCCESS explicitly targets the interaction between market structures, long-term contracting, and public policy. This is also where my strengths are most useful. I have experience building and stress-testing models where contracts or rules discipline behavior over time (e.g., investment dynamics, price adjustments, and state-dependent strategies), and where regulation changes feasible actions and equilibrium outcomes. I am especially interested in research questions such as: Which contracting structures best mitigate investment hold-up? How do long-term commitments affect entry and competition? How should market rules balance risk-sharing, liquidity, and dynamic efficiency during the scale-up of a new energy market?

I am also motivated by the project’s multidisciplinary setup: collaboration with partners at TU Delft and the integration of legal-institutional perspectives through TILEC. My work is designed to be structurally interpretable and policy-relevant---a style that translates naturally to multidisciplinary work where models must connect to engineering constraints (e.g., capacity, transport bottlenecks, storage) and institutional feasibility (e.g., governance, regulation, contracting enforceability). In terms of tools, I work primarily in Python and R and have extensive experience producing reproducible computational pipelines for simulation, estimation, and policy evaluation.

Beyond my dissertation, my coauthored projects demonstrate additional strengths useful for HySUCCESS: (i) \emph{Polytope Fraud Theory} applies machine learning to large-scale financial data to detect structured anomalies, reflecting strong quantitative modeling and data-handling skills; (ii) \emph{Agentic Workflows for Economic Research} develops automation pipelines for empirical work, reflecting my ability to build scalable and transparent research infrastructure. These experiences are transferable to a large collaborative project environment that involves data integration, scenario analysis, and communication of results to academic and policy audiences.

If selected, I would aim to contribute to HySUCCESS by:
\begin{itemize}
  \item developing economic models of hydrogen market structure and regulation, with explicit treatment of dynamic investment incentives;
  \item analyzing the design and welfare implications of long-term contracting arrangements (with links to contract theory and market design);
  \item conducting simulation-based counterfactuals to compare alternative rule and contract designs under uncertainty;
  \item supporting interdisciplinary translation between economic modeling, engineering constraints, and legal-institutional requirements, in collaboration with TU Delft and TILEC partners.
\end{itemize}

Thank you for considering my application. I would be delighted to contribute to Tilburg University’s research community and to the HySUCCESS / GroenvermogenNL programme under the supervision of Dr.\ Bert Willems and in collaboration with TILEC and TU Delft partners. I would welcome the opportunity to discuss how my background in industrial organization and computational modeling can support the development of robust and economically sustainable hydrogen market designs.

\vspace{0.5cm}

Sincerely,\\[0.8cm]
Zhongli Wang

\end{document}
